\documentclass[10pt,a4paper]{article}
\usepackage[spanish,es-nodecimaldot]{babel}
\usepackage[utf8]{inputenc}
\usepackage[T1]{fontenc}
\usepackage{lmodern}
\usepackage[top=1.1cm,bottom=1.1cm,left=1.5cm,right=1.5cm]{geometry}
\usepackage{amsmath,amssymb}
\usepackage[table]{xcolor}
\usepackage{booktabs,array}
\usepackage{enumitem}
\usepackage[normalem]{ulem}

\setlength{\parindent}{0pt}
\setlength{\parskip}{2pt}
\pagestyle{empty}

\AtBeginDocument{%
  \setlength{\abovedisplayskip}{3pt}%
  \setlength{\belowdisplayskip}{3pt}%
  \setlength{\abovedisplayshortskip}{2pt}%
  \setlength{\belowdisplayshortskip}{2pt}%
}

\begin{document}

{\large\bfseries Diseño en bloques aleatorios}
\vspace{2pt}

El experimentador tiene a disposicion mediciones relativas a ``a'' tratamientos distribuidos en b bloques. Vamos a considerar el caso en q hay una observacion en cada tratamiento de cada bloque.

\begin{center}
\vspace{-2pt}
{\small
\begin{tabular}{l c c c c | c c c}
\toprule
Tratamiento & Bloque 1 & Bloque 2 & $\ldots$ & Bloque $b$ & Total ($T_{i.}$) & Media ($\bar{y}_{i.}$) & Efecto ($\alpha_i$) \\
\midrule
Tratamiento 1 & $y_{11}$ & $y_{12}$ & $\ldots$ & $y_{1b}$ & $T_{1.}$ & $\bar{y}_{1.}$ & $\alpha_1=\bar{y}_{1.}-\bar{y}_{..}$ \\
Tratamiento 2 & $y_{21}$ & $y_{22}$ & $\ldots$ & $y_{2b}$ & $T_{2.}$ & $\bar{y}_{2.}$ & $\alpha_2=\bar{y}_{2.}-\bar{y}_{..}$ \\
$\ldots$ & $\ldots$ & $\ldots$ & $\ldots$ & $\ldots$ & $\ldots$ & $\ldots$ & $\ldots$ \\
Tratamiento $a$ & $y_{a1}$ & $y_{a2}$ & $\ldots$ & $y_{ab}$ & $T_{a.}$ & $\bar{y}_{a.}$ & $\alpha_a=\bar{y}_{a.}-\bar{y}_{..}$ \\
\midrule
Total Bloque & $T_{.1}$ & $T_{.2}$ & $\ldots$ & $T_{.b}$ & $T_{..}$ & & \\
Media Bloque & $\bar{y}_{.1}$ & $\bar{y}_{.2}$ & $\ldots$ & $\bar{y}_{.b}$ & & $\bar{y}_{..}$ & \\
Efecto Bloque & $\beta_1$ & $\beta_2$ & $\ldots$ & $\beta_b$ & & & \\
\bottomrule
\end{tabular}
}
\end{center}

\hfill {\small a=numero de tratamiento \quad b=numero de bloque}

El modelo matemático aditivo sin interacción es:
\[
y_{ij} = \mu + \alpha_i + \beta_j + \epsilon_{ij} \qquad \text{para } i=1,\ldots,a; \quad j=1,\ldots,b
\]

Con las restricciones $\sum_{i=1}^{a}\alpha_i=0$ y $\sum_{j=1}^{b}\beta_j=0$, donde:
\begin{itemize}[leftmargin=16pt,itemsep=0pt,topsep=1pt,parsep=0pt]
  \item $\mu$: Gran media general.
  \item $\alpha_i=\bar{y}_{i.}-\bar{y}_{..}$: Efecto fijo del $i$-ésimo tratamiento.
  \item $\beta_j=\bar{y}_{.j}-\bar{y}_{..}$: Efecto fijo del $j$-ésimo bloque.
  \item $\epsilon_{ij} \sim iid\ N(0,\sigma^2)$: Error aleatorio.
\end{itemize}

\hfill {\small $\alpha_i$ =efecto de la fila \quad $\beta_j$= efecto de la columna \quad $\varepsilon_{ij}$ = ruido o error que tiene media cero o varianza comun $\sigma^2$}

\[
C = \frac{T_{..}^2}{a\cdot b}
\]
\[
SST = \sum_{i=1}^{a}\sum_{j=1}^{b} y_{ij}^2 - C
\]
\[
SS(Tr) = \frac{1}{b}\sum_{i=1}^{a} T_{i.}^2 - C \qquad \text{con } a-1 \text{ gl}
\]
\[
SS(Bl) = \frac{1}{a}\sum_{j=1}^{b} T_{.j}^2 - C \qquad \text{con } b-1 \text{ gl}
\]
\[
SSE = SST - SS(Tr) - SS(Bl) \qquad \text{con } (a-1)(b-1) \text{ gl}
\]

\begin{center}
\vspace{-2pt}
{\footnotesize
\begin{tabular}{l c c c c}
\toprule
\textbf{Fuentes de Variación} & \textbf{Grados de Libertad} & \textbf{Suma de Cuadrados} & \textbf{Media Cuadrada} & \textbf{F} \\
\midrule
Tratamientos & $a-1$ & $SS(Tr)$ & $MS(Tr)=SS(Tr)/(a-1)$ & $F_{Tr}=MS(Tr)/MSE$ \\
Bloques & $b-1$ & $SS(Bl)$ & $MS(Bl)=SS(Bl)/(b-1)$ & $F_{Bl}=MS(Bl)/MSE$ \\
Error & $(a-1)\cdot(b-1)$ & $SSE$ & $MSE=SSE/(a-1)\cdot(b-1)$ & \\
Total & $a\cdot b-1$ & $SST$ & & \\
\bottomrule
\end{tabular}
}
\end{center}

para llegar al error hicimos: \quad Error = total-tratamientos-bloques

\end{document}
